\documentclass[10pt,a4paper]{amsart}
\usepackage[margin=19mm]{geometry}
\usepackage{amsmath,amssymb,mathtools}
\usepackage[scr=rsfs,cal=euler]{mathalfa}
\usepackage{xfrac}
\usepackage[unicode,hidelinks,bookmarks=false]{hyperref}
\newcommand{\ie}{i.e.\ }
\newcommand{\cf}{cf.\ }
\newcommand{\resp}{respectively\ }
\newcommand{\Subsection}{\subsection}
\providecommand{\nobreakdash}{\nobreak\hbox{-}}
\setlength{\emergencystretch}{2em}
\multlinegap0pt
\usepackage{bm}
\usepackage{bbm}
\usepackage{dsfont}
\usepackage{xspace}
\usepackage{cancel}
\usepackage{mathtools}
\usepackage{amsthm}
\makeatletter
\@ifundefined{lemma}{\newtheorem{lemma}{Lemma}}{}
\@ifundefined{theorem}{\newtheorem{theorem}[lemma]{Theorem}}{}
\makeatother
\newcommand\be{\begin{equation}}
\newcommand\ee{\end{equation}}
\newcommand{\B}{\mathbb{B}}
\newcommand{\Values}{worths}
\newcommand{\Value}{worth}
\title[Emergence in graphs with near-extreme constraints]{Emergence in graphs with near-extreme constraints}
\author{Charles Radin, Lorenzo Sadun}
\date{Source: arXiv:2411.14556v3 (CC BY 4.0)}
\begin{document}
\maketitle

\noindent\emph{Source and licence.} Emergence in graphs with near-extreme constraints. Version arXiv:2411.14556v3 (CC BY 4.0), distributed under the Creative Commons Attribution 4.0 International licence, as recorded in the arXiv licence metadata for this exact version and shown on the abstract page https://arxiv.org/abs/2411.14556v3. Full rights notice: licenses/arxiv-2411.14556-CC-BY-4.0.txt.\par\smallskip

\noindent\textit{Editorial scope.} Counted unit: the Theorem whose source label is \texttt{thm:top2} in the article (source0.tex lines 1814--1819, source label \texttt{thm:top2}), delivered together with the article's own complete original proof of it (source0.tex lines 1821--1919, one \texttt{proof} environment of 99 lines). No mathematics is written, repaired or re-derived: statement and proof are the article's own text line for line. Required context retained uncounted: thm:no-singular (lines 897--901); thm:flat2 (lines 930--936); thm:scallop2 (lines 1326--1338). Every definition, display and label of those lines is retained unchanged. Omitted from this package and not redistributed: 1--896, 902--929, 937--1325, 1339--1813, 1820--1820, 1920--2285. Nothing inside a retained excerpt is omitted or altered: every retained body is delivered exactly as the article's lines are written, so no omission receipt of any kind accompanies this package. Citations: the retained excerpts carry no citation command at all, so no bibliography entry of the article is transcribed and no reference is redistributed. Reference adaptation: none was needed and none was made; every reference inside the delivered unit resolves inside it, so no unresolved reference or citation is produced. Printed slips kept literal: no printed word, symbol, hypothesis or number of the counted statement or of its proof was altered. Layout and class: the article's own class is \texttt{amsart} with packages including amsfonts, amssymb, amsmath, graphicx, dcolumn, bm, xcolor, comment, ulem, subfigure; the wrapper is the lane's pinned \texttt{amsart} editorial wrapper at 10pt on A4, which restates the theorem-like environments the retained text needs and adds the article's own macro definitions that the retained text needs (\texttt{\textbackslash B}, \texttt{\textbackslash Value}, \texttt{\textbackslash Values}, \texttt{\textbackslash be}, \texttt{\textbackslash ee}), with extra wrapper packages (bm, bbm, dsfont, xspace, cancel, mathtools). The delivered numbering is the unit's own. Assets: no retained span contains a figure environment, a graphics-inclusion command, a PDF-page inclusion, a table or a TikZ picture; no image file is copied into this package, the package declares no asset dependency and no image file is redistributed with it. Licence: version arXiv:2411.14556v3 of this article is distributed under the Creative Commons Attribution 4.0 International licence, as recorded in the arXiv licence metadata for this exact version and shown on the abstract page https://arxiv.org/abs/2411.14556v3. A full rights notice is delivered at licenses/arxiv-2411.14556-CC-BY-4.0.txt. Characters: every retained excerpt is pure ASCII, so the delivered engine, XeLaTeX with its default Latin Modern text font, reports no missing character.
\begin{theorem} \label{thm:no-singular} 
Fix the edge density $e$ and consider an open interval $I$ of triangle densities $t$. Suppose that there is a
differentiable function $f(t)$, defined for all $t \in I$, that equals the Boltzmann entropy $\B(e,t)$ 
whenever $t$ is non-singular. Then every $t \in I$ is non-singular. 
\end{theorem} 

\begin{theorem} \label{thm:flat2}
For each fixed $e<1/2$ and all $t$ sufficiently small, the 
optimal graphon with edge/triangle densities $(e,t)$ is unique and is symmetric bipodal, with parameters that vary analytically with $(e,t)$. As $t \to 0$,
the increase $\Delta \B = \B(e,t)-\B(e,0)$ in the 
Boltzmann entropy scales as $t \ln(1/t)$ and the Lagrange multiplier 
$\beta$ scales as $\ln(1/t)$. 
\end{theorem} 

\begin{theorem} \label{thm:scallop2}
Let $n \ge 1$ be an integer. For every $e \in \left ( \frac{n}{n+1}, \frac{n+1}{n+2}\right )$, with 
corresponding minimal triangle density $t_0$ (depending on $e$), 
and for all
$\Delta t$ sufficiently small, the optimal graphon with edge/triangle densities $(e, t_0 + \Delta t)$
is unique and $n+2$-podal, with $(n,2)$ symmetry. Asymptotically, 
$\Delta \B = \B(e,t) - 
\B(e,t_0)$ scales as $\sqrt{\Delta t}$.
and the Lagrange multiplier 
$\beta$ scales as $1/\sqrt{\Delta t}$.
In the optimal graphon, the diagonal entries are all $\exp(-\Theta(\beta))$ and, except for the 
$(n+1, n+2)$ and $(n+2,n+1)$ entries, the off-diagonal entries are all $1 - \exp(-\Theta(\beta))$.
\end{theorem}

\begin{theorem}\label{thm:top2}
For each fixed $e \in (0,1)$ and all $t$ sufficiently close to (but below) $e^{3/2}$, the 
optimal graphon with edge/triangle densities $(e,t)$ is unique and bipodal. Asymptotically, the Boltzmann entropy 
scales as $-(e^{3/2}-t)\ln(e^{3/2}-t)$
and the Lagrange multiplier $\beta$ scales as $\ln(e^{3/2}-t)$. 
\end{theorem}

\begin{proof}
We follow the same overall roadmap as 
the proofs of Theorems \ref{thm:flat2} and \ref{thm:scallop2}. 
Specifically, 
\begin{enumerate}
\item Using the proximity to the upper boundary, we break $[0,1]$
into two large podes $I_1$ and $I_2$ and a small exceptional set 
$I_3$ such 
that $g_x$ is $L^2$-close to the indicator function of $I_1$ when
$x \in I_1$ and is $L^2$-close to zero when $x \in I_2$. 
\item Equating the \Values{} of $g_x$ when $x \in I_1$ to those of 
$g_x$ when $x \in I_2$, we determine that $\beta/\alpha
\approx -2/e$. The multiplier $\beta$ is large and negative,
while $\alpha$ is large and positive. 
\item Maximizing $W(C)$ for an arbitrary $C: [0,1] \to [0,1]$, we show 
that every column is close to a typical column in the first or second
pode. After reassigning points,
$I_3$ is then empty. The control this gives us on $G(x,y)$ shows that 
$g(x,y)$ is everywhere exponentially close to 0 or 1.
\item Bounding the fluctuations in each rectangle by multiples 
of the fluctuations in other rectangles to show that all 
fluctuations are in fact zero. In other words, our optimal graphon is 
exactly bipodal with values that are exponentially close to 0 or 1. 
\item Using Theorem \ref{thm:no-singular} to eliminate the possibility
that some points $(e,t)$ might have entropy maximizers without well-defined Lagrange multipliers 
$(\alpha, \beta)$.
\end{enumerate}

Step 1 is identical to what we have done before. There is a 
unique reduced graphon at $(e, e^{3/2})$,
namely the equivalence class of a graphon $g_0$ that is 1 on $I_1 \times I_1$ and zero elsewhere, where $I_1$ is a pode of size $\sqrt{e}$. Every graphon with $t$ 
close to $e^{3/2}$, and in particular any entropy-maximizing graphon,
must be $L^2$ close to $g_0$. This means that for all $x$'s outside
of a set of small measure, $g_x$ is $L^2$-close to the corresponding
column of $g_0$. This also implies that $G(x,y)$ is close to $\sqrt{e}$
when $x$ and $y$ are both in $I_1$ and is close to zero when either
is in $I_2$. 

The \Value{} of a column that is nearly zero is of course nearly zero.
The \Value{} of a column that is nearly 1 on $I_1$ and nearly zero
elsewhere is approximately 
\be - 2\alpha \sqrt{e} - \beta e^{3/2}. \ee
Since all columns must have the same \Value, we must have 
$\beta/\alpha \approx -2/e$. The Lagrange multipliers $\alpha$
and $\beta$ diverge at the same rate as we approach the boundary,
with $\alpha \to \infty$ and $\beta \to -\infty$. 

Now consider an arbitrary function $C: [0,1] \to [0,1]$. 
Let $a$ be the 
average of $C(y)$ on $I_1$ and let $b$ be the average on $I_2$. 
Using the approximation that $G$ is $L^2$-close to $\sqrt{e}$ 
times the indicator function of $I_1 \times I_1$, we get that 
\be W(C) \le 2\sqrt{e} H(a) + 2(1-\sqrt{e}) H(b) - 2\alpha \sqrt{e} a
- 2\alpha (1-\sqrt{e}) b - \beta e^{3/2} a^2, \ee
with equality if $C$ is constant on $I_1$ and constant on $I_2$. 
Since $\alpha$ is large and positive, we must have $b$ exponentially
close to 0. Setting $b \approx 0$, our \Value{} is then approximately
\be \alpha \sqrt{e} (a^2 - a). \ee
This is of course maximized at the endpoints $a=1$ and $a=0$, 
being negative when $a \in (0,1)$. In other words, 
any \Value-maximizing column must either have $a \approx 1$ and 
$b \approx 0$, and so must be close to the columns in $I_1$, or 
$a \approx 0$ and $b \approx 0$, and so must be close to the 
columns in $I_2$. Reassigning the points of $I_3$ to $I_1$ or $I_2$
accordingly, we obtain a situation where $I_3$ is empty. 

To constrain the fluctuations in $g(x,y)$ in each rectangle, we
recall the variational equations 
\be H'(g(x,y)) = \alpha + \beta G(x,y). \ee
Since $G(x,y) \approx 0$ or $\sqrt{e}$, depending on which 
quadrant we are in, this implies that $g(x,y)$ is exponentially 
close to 1 on $I_1 \times I_1$ and exponentially close to 
0 on $I_1 \times I_2$ and $I_2 \times I_2$. In particular, $H''(g)$,
which scales as the larger of $1/g$ and $1/(1-g)$, is much larger
than $|\beta|$. Looking at the change in the left hand side and 
right hand side of this equation within a single quadrant, we 
see that $H''(g)$ times the maximum fluctuation within any quadrant
is of the same order as $|\beta|$ times the maximum fluctuation 
within any quadrant. But that means that the maximum fluctuation is 
bounded by a small multiple of itself, and so must be zero. Our 
graphon is exactly bipodal. 

Finally, we do some calculations in the space of bipodal graphons. 
Let $g_{11}$, $g_{12}$ and $g_{22}$ be the values of the
optimal graphon on $I_1 \times I_1$, $I_1 \times I_2$ and $I_2 \times I_2$,
respectively. We treat $g_{11}$, $g_{12}$ and $g_{22}$ as free variables
and adjust the width of $I_1$ to keep the edge density fixed. To 
leading order, $e^{3/2}-t$ is a linear function of $(1-g_{11}, g_{12},
g_{22})$. However, $1-g_{11}$, $g_{12}$ and $g_{22}$ all scale as 
exponents of $\beta$, so $\beta$ must scale as $\ln(e^{3/2}-t)$. 
The entropy goes as $-(1-g_{11})\ln(1-g_{11}) - g_{12} \ln(g_{12})
- g_{22} \ln(g_{22})$, which then scales as $-(e^{3/2}-t)) \ln(e^{3/2}-t)$.

The analyticity of $g$ as a function of $(e,t)$ follows from the 
same argument as in the proof of Theorem \ref{thm:scallop2}, only with
the space of $(n,2)$ symmetric graphons replaced by the space of
bipodal graphons. 

\end{proof}

\end{document}
